% Options for packages loaded elsewhere
\PassOptionsToPackage{unicode}{hyperref}
\PassOptionsToPackage{hyphens}{url}
\PassOptionsToPackage{dvipsnames,svgnames,x11names}{xcolor}
%
\documentclass[
  11pt,
  a4paper,
]{ltjsarticle}
\usepackage{amsmath,amssymb}
\usepackage{iftex}
\ifPDFTeX
  \usepackage[T1]{fontenc}
  \usepackage[utf8]{inputenc}
  \usepackage{textcomp} % provide euro and other symbols
\else % if luatex or xetex
  \usepackage{unicode-math} % this also loads fontspec
  \defaultfontfeatures{Scale=MatchLowercase}
  \defaultfontfeatures[\rmfamily]{Ligatures=TeX,Scale=1}
\fi
\usepackage{lmodern}
\ifPDFTeX\else
  % xetex/luatex font selection
\fi
% Use upquote if available, for straight quotes in verbatim environments
\IfFileExists{upquote.sty}{\usepackage{upquote}}{}
\IfFileExists{microtype.sty}{% use microtype if available
  \usepackage[]{microtype}
  \UseMicrotypeSet[protrusion]{basicmath} % disable protrusion for tt fonts
}{}
\makeatletter
\@ifundefined{KOMAClassName}{% if non-KOMA class
  \IfFileExists{parskip.sty}{%
    \usepackage{parskip}
  }{% else
    \setlength{\parindent}{0pt}
    \setlength{\parskip}{6pt plus 2pt minus 1pt}}
}{% if KOMA class
  \KOMAoptions{parskip=half}}
\makeatother
\usepackage{xcolor}
\usepackage[margin=24mm]{geometry}
\setlength{\emergencystretch}{3em} % prevent overfull lines
\providecommand{\tightlist}{%
  \setlength{\itemsep}{0pt}\setlength{\parskip}{0pt}}
\setcounter{secnumdepth}{5}
\ifLuaTeX
  \usepackage{selnolig}  % disable illegal ligatures
\fi
\IfFileExists{bookmark.sty}{\usepackage{bookmark}}{\usepackage{hyperref}}
\IfFileExists{xurl.sty}{\usepackage{xurl}}{} % add URL line breaks if available
\urlstyle{same}
\hypersetup{
  colorlinks=true,
  linkcolor={blue},
  filecolor={Maroon},
  citecolor={Blue},
  urlcolor={blue},
  pdfcreator={LaTeX via pandoc}}

\author{}
\date{}

\begin{document}

\hypertarget{chapter-03-ux9023ux7acbux4e00ux6b21ux65b9ux7a0bux5f0f}{%
\section{Chapter 03 ---
連立一次方程式}\label{chapter-03-ux9023ux7acbux4e00ux6b21ux65b9ux7a0bux5f0f}}

\texttt{Ax=b} は、単に未知数を求める問題ではありません。

\hypertarget{ux4e09ux3064ux306eux540cux3058ux554fux984c}{%
\subsection{三つの同じ問題}\label{ux4e09ux3064ux306eux540cux3058ux554fux984c}}

\[
Ax=b
\]

は次の三つを同時に表します。

\begin{itemize}
\tightlist
\item
  方程式をすべて満たす \texttt{x} を探す。
\item
  \texttt{b} を \texttt{A} の列の一次結合で作る。
\item
  線形写像 \texttt{A} で \texttt{b} へ移る入力を探す。
\end{itemize}

このため「解が存在するか」は

\[
b\in Col(A)
\]

という部分空間の問題です。

\hypertarget{ux81eaux7531ux5909ux6570ux306eux672cux5f53ux306eux610fux5473}{%
\subsection{自由変数の本当の意味}\label{ux81eaux7531ux5909ux6570ux306eux672cux5f53ux306eux610fux5473}}

一つの解 \texttt{x\_p} があるとき

\[
x=x_p+z,\qquad z\in\ker A
\]

です。

自由変数は計算途中に偶然現れるものではなく、\texttt{A}
が見分けられない零空間方向の自由度です。

\hypertarget{gauss-ux6d88ux53bbux3092ux3069ux3046ux7406ux89e3ux3059ux308bux304b}{%
\subsection{Gauss
消去をどう理解するか}\label{gauss-ux6d88ux53bbux3092ux3069ux3046ux7406ux89e3ux3059ux308bux304b}}

消去法は解を出すアルゴリズムであると同時に、

\begin{itemize}
\tightlist
\item
  pivot = 独立な制約
\item
  free variable = 核方向
\item
  inconsistent row = \texttt{b} が列空間外
\end{itemize}

を露出させる方法です。

\hypertarget{ux5230ux9054ux57faux6e96}{%
\subsection{到達基準}\label{ux5230ux9054ux57faux6e96}}

Gauss 消去の結果を見て、解の個数だけでなく
rank、nullity、解空間の幾何学まで説明できることを目標にしてください。

\end{document}
